 \chapter{Example}
\label{chp:example} 

Here is an example of how to use acronyms such as \gls{ntnu}. The second time only \gls{ntnu} is shown and if there were several you would write \glspl{ntnu}. And here is an example\footnote{A footnote} of citation~\cite{Author:year:XYZ}. 

\Blindtext[3][1]

\begin{figure}
\centering
% dummy figure replacement 
\begin{tabular}{@{}c@{}}
\rule{.5\textwidth}{.5\textwidth} \\
\end{tabular}
\caption{\label{fig:example}A figure}
\end{figure}

% This shows how to deal with titles that are too wide for the header, see https://www.texfaq.org/FAQ-runheadtoobig for details
% For short titles \section{First section} is enough
\section[First section (middling version)(ToC)]{First section (verbose version)%
              \sectionmark{First section (terse version)(head current)}}
\sectionmark{First section (terse version)(head next)}
\label{sec:first_section}

\subsection{First subsection with some \texorpdfstring{$\mathcal{M}ath$}{Math} symbol}\label{sec:first_ssection}

\blindtext
\begin{itemize}
 \item item1
 \item item2
 \item ...
\end{itemize}

\subsection{Mathematics}

\blindmathtrue
\blindtext

\begin{proposition}\label{def:a_proposition}
A proposition... (similar environments include: theorem, corollary, conjecture, lemma)

\end{proposition}

\begin{proof}
\vspace*{-1em} % Adjust the space when parskip is set to 1em
And its proof.
\end{proof}

\begin{table}
\caption{\label{tab:example}A table}
\centering
\begin{tabular}[b]{| c | c | c | c | c |}
\hline
a & b & c & d & e \\ \hline
f & g & h & i & j \\ \hline
k & l & m & n & o \\ \hline
p & q & r & s & t \\ \hline
u & v & w & x & y \\ \hline
z & æ & ø & å &   \\ \hline
\end{tabular} 
\end{table}

\subsection{Source code example}

% \floatname{algorithm}{Source code} % if you want to rename 'Algorithm' to 'Source code'
\begin{algorithm}[h]
  \caption{The Hello World! program in Java.}
  \label{hello_world}
  % alternatively you may use algorithmic, or lstlisting from the listings package
  \begin{verbatim}
  
class HelloWorldApp {
  public static void main(String[] args) {
    //Display the string
    System.out.println("Hello World!");
  }
}
\end{verbatim}
\end{algorithm}

You can refer to figures using the predefined command like \fref{fig:example}, to pages like \pref{fig:example}, to tables like \tref{tab:example}, to chapters like \Cref{chp:example} and to sections like \Sref{sec:first_section} and you may define similar commands to refer to proposition, algorithms etc.
